\clearpage
\section*{ЗАКЛЮЧЕНИЕ}
\addcontentsline{toc}{section}{ЗАКЛЮЧЕНИЕ}

В ходе выполнения курсового проекта была достигнута поставленная цель --- разработано и обосновано программное решение \textit{Gamified Task CRM (Flutter Web)}, ориентированное на управление задачами и клиентскими взаимодействиями в едином цифровом контуре. Работа показала, что даже в рамках учебного проекта возможно реализовать систему, в которой объединены предметная логика CRM, структурированная архитектура, backend-first подход к данным и формализованный контроль качества.

Первая задача, связанная с анализом теоретических и технологических основ разработки CRM-ориентированного приложения, выполнена в полном объеме. В теоретической части рассмотрены назначение CRM-систем, требования к функциональному и нефункциональному контуру, подходы к архитектуре, принципы проектирования данных и базовые практики надежности и безопасности. На этой основе обоснован выбор стека Flutter + Riverpod + GoRouter + Supabase как рационального решения для проекта с требованиями к модульности, прозрачной маршрутизации и реляционной модели данных.

Вторая задача, предусматривающая проектирование логической структуры решения, также выполнена. В проекте выделены функциональные модули интерфейса (dashboard, tasks, crm, reports, profile, ai), определены границы ответственности, реализована единая оболочка навигации и согласованный шаблон пользовательских действий. На уровне данных зафиксирована ключевая связь \textit{задача-клиент} через \texttt{client\_id}, что обеспечивает трассируемость операционной активности. Подобная структура позволяет последовательно масштабировать систему без радикального пересмотра архитектуры.

Третья задача, связанная с практической реализацией, решена путем построения рабочего web-приложения на Flutter с поддержкой авторизации, управления клиентами и задачами, базовой аналитики и режимов исполнения. Реализация содержит как backend-first контур для реального использования, так и demo-режим для воспроизводимой учебной демонстрации. Такой подход позволил совместить практичность и устойчивость: система работает с внешним backend, но при необходимости может быть показана в автономном режиме без потери основной функциональности.

Четвертая задача --- проверка корректности и устойчивости решения --- выполнена через многоуровневый quality-контур: статический анализ, автоматизированные тесты, smoke-проверки и web E2E-сценарий. В результате подтверждена работоспособность ключевого пользовательского процесса \textit{Login $\rightarrow$ Create Client $\rightarrow$ Create Task $\rightarrow$ Verify}. Применение формализованных проверок повысило доверие к результатам разработки и обеспечило воспроизводимость демонстрации перед защитой.

Итоговые результаты курсового проекта можно свести к следующим положениям:
\begin{enumerate}
  \item Создано прикладное CRM-ориентированное приложение с модульной архитектурой и четкой декомпозицией функционала.
  \item Обеспечена связность операционных сущностей, включая корректную модель данных для задач, клиентов и пользовательской статистики.
  \item Реализован безопасный конфигурационный подход через переменные окружения и разграничение demo/production режимов.
  \item Сформирован практический контур верификации качества, достаточный для учебной и предэксплуатационной проверки.
  \item Подготовлена база для дальнейшего развития проекта без нарушения базовых архитектурных принципов.
\end{enumerate}

Практическая значимость работы заключается в том, что разработанное решение может использоваться как учебный стенд полного цикла разработки программного обеспечения. Проект позволяет демонстрировать не только интерфейсную часть, но и инженерные дисциплины: организацию данных, контроль качества, управление конфигурацией, обработку ошибок, подготовку к деплою. Это делает работу применимой как для образовательных задач, так и для стартовых прикладных сценариев малого бизнеса или учебных лабораторий.

Одновременно следует обозначить ограничения текущей версии. В системе пока отсутствуют расширенные аналитические инструменты уровня BI, полноценная ролевая модель с многоуровневым разграничением прав и глубокий нагрузочный контур. Кроме того, AI-модуль играет вспомогательную роль и требует дополнительного исследования в части качества ответов и прикладной полезности для CRM-процессов. Тем не менее эти ограничения являются ожидаемыми для масштаба курсового проекта и не противоречат достижению его основной цели.

С учетом полученных результатов можно предложить следующие направления дальнейшего развития:
\begin{enumerate}
  \item Расширить отчетный блок: добавить воронку продаж, периодическую динамику и сегментацию клиентской базы.
  \item Ввести ролевую модель доступа и аудит критичных операций.
  \item Реализовать сервис уведомлений по дедлайнам задач и событиям CRM.
  \item Усилить качество через дополнительные интеграционные и контрактные тесты.
  \item Масштабировать продукт в сторону mobile-first сценариев на общей Flutter-кодовой базе.
\end{enumerate}

Подводя итог, можно заключить, что курсовой проект демонстрирует достижение поставленной цели и решение всех сформулированных задач. Разработанное программное обеспечение обладает практической применимостью, соответствует требованиям к структуре и содержанию учебной работы и служит устойчивой основой для последующего инженерного развития.
